\begin{pastexam}[label={px:nccu-106-dsa-06}]{政大資科所 106 學年度 第 6 題}
（問答）

Recall the traveling salesman problem (TSP):

\hspace*{1.5em}Input: A list of cities; the distances between each pair of cities; a budget $b$.\\
\hspace*{1.5em}Output: A route of length less than or equal to $b$ that visits each city exactly
once and returns to the origin city, if such a route exists.

The optimization version of this problem asks directly for the shortest route (TSP-OPT):

\hspace*{1.5em}Input: A list of cities; the distances between each pair of cities.\\
\hspace*{1.5em}Output: The shortest possible route that visits each city exactly once and returns
to the origin city.
\begin{enumerate}[label=\Alph*.]
\item (5\%) Show that if TSP-OPT can be solved in polynomial time, then so can TSP.
\item (5\%) Show that if TSP can be solved in polynomial time, then so can TSP-OPT.
\end{enumerate}
\end{pastexam}

\begin{pastsoln}
\textbf{約定}。 城市有 $n$ 個，距離是非負整數，最大的距離是 $C$（距離是有理數時，全部乘上分母的
公倍數就變成整數，不影響哪一條路線最短）。 輸入長度至少要寫下每個距離，所以 $n$ 和 $\lg C$ 都不超過
輸入長度。 題目的 TSP 是\textbf{搜尋版}：有符合預算的路線就輸出一條，沒有就回答沒有。 假設 TSP-OPT
的演算法是 $\mathcal{A}_{\text{OPT}}$，TSP 的演算法是 $\mathcal{A}$，都是多項式時間。

\textbf{A. TSP-OPT 好解 $\Rightarrow$ TSP 好解}。 給 TSP 的輸入（距離和預算 $b$）：
\begin{enumerate}
\item 用 $\mathcal{A}_{\text{OPT}}$ 求出最短路線 $R$。
\item 花 $O(n)$ 把 $R$ 的 $n$ 段距離加起來，得到長度 $L$。
\item 若 $L \le b$，輸出 $R$；否則回答「沒有這樣的路線」。
\end{enumerate}
正確性：如果有一條長度 $\le b$ 的路線，最短的路線長度 $L$ 不會更長，所以 $L \le b$，輸出的 $R$
符合要求；如果 $L > b$，所有路線都至少長 $L$，都超過預算，回答沒有是對的。 時間是一次
$\mathcal{A}_{\text{OPT}}$ 加上 $O(n)$ 個加法，多項式時間。

\textbf{B. TSP 好解 $\Rightarrow$ TSP-OPT 好解}。 最短路線的長度 $\mathrm{OPT}$ 是整數，而且
$0 \le \mathrm{OPT} \le nC$（路線有 $n$ 段，每段 $\le C$）。 $\mathcal{A}$ 在預算 $m$ 時找得到路線，
若且唯若 $\mathrm{OPT} \le m$。 所以可以對 $m$ 做二分搜尋（資結書第~28.1.3~節）：
\begin{pseudo}[unbreakable]{用 TSP 解 TSP-OPT}
\Function{TSP-Opt}{$D$} \Comment{$D$：城市間的距離}
    \State $\mathit{lo} \gets 0$；$\mathit{hi} \gets nC$ \Comment{不變量：$\mathit{lo} \le \mathrm{OPT} \le \mathit{hi}$}
    \While{$\mathit{lo} < \mathit{hi}$}
        \State $m \gets \lfloor (\mathit{lo} + \mathit{hi}) / 2 \rfloor$
        \If{$\mathcal{A}(D, m)$ 找到路線}
            \State $\mathit{hi} \gets m$ \Comment{$\mathrm{OPT} \le m$}
        \Else
            \State $\mathit{lo} \gets m + 1$ \Comment{$\mathrm{OPT} > m$}
        \EndIf
    \EndWhile
    \State \Return $\mathcal{A}(D, \mathit{lo})$ 找到的路線
\EndFunction
\end{pseudo}

正確性：一開始 $0 \le \mathrm{OPT} \le nC$，不變量成立。 第 5 行找到路線表示有長度 $\le m$ 的路線，
$\mathrm{OPT} \le m$；找不到表示每條路線都 $> m$，$\mathrm{OPT} \ge m + 1$。 兩種更新都維持不變量，
而且區間每輪至少縮小一個，迴圈會結束。 結束時 $\mathit{lo} = \mathit{hi} = \mathrm{OPT}$。 第 9 行在
預算 $\mathrm{OPT}$ 下一定找得到路線，它的長度 $\le \mathrm{OPT}$，而沒有路線比 $\mathrm{OPT}$ 短，
所以它的長度恰好是 $\mathrm{OPT}$，就是一條最短路線。

時間：區間長度從 $nC + 1$ 每輪減半，迴圈最多跑 $\lceil \lg(nC + 1) \rceil = O(\lg n + \lg C)$ 輪，
這是輸入長度的多項式（其實是線性）。 每次呼叫 $\mathcal{A}$ 的輸入是原本的距離加上一個不超過 $nC$ 的
預算，預算只多 $O(\lg n + \lg C)$ 個位元，輸入長度還是原來的多項式，所以每次呼叫是多項式時間。
呼叫 $O(\lg n + \lg C)$ 次 $\mathcal{A}$，總時間是多項式。

注意不能從 $m = 0, 1, 2, \ldots$ 一個一個試：最多要試 $nC + 1$ 次，$C$ 可以是輸入長度的指數
（$C$ 只要 $\lg C$ 個位元就寫得下），那樣不是多項式時間。

\textbf{補充}。 如果 TSP 只回答「有／沒有」（決策版），還要把路線找出來：求出 $\mathrm{OPT}$ 之後，
依序把每一對城市的距離改成 $\mathrm{OPT} + 1$（等於禁止走這段），改了之後預算 $\mathrm{OPT}$ 仍然回答
「有」就保留這個修改，否則改回來。 全部試完，沒被禁止的那 $n$ 段就是一條最短路線（資結書第~28.1.3~節的
自我歸約），總共再多 $\binom{n}{2}$ 次呼叫，仍是多項式時間。
\end{pastsoln}
